\documentclass[11pt,a4paper]{article}
\usepackage[utf8]{inputenc}
\usepackage[T1]{fontenc}
\usepackage[ngerman]{babel}
\usepackage{geometry}
\geometry{margin=2.5cm}
\usepackage{booktabs}
\usepackage{enumitem}
\usepackage{xcolor}
\usepackage{scrextend}
\usepackage{hyperref}
\usepackage{tabularx}
\usepackage{longtable}
\hypersetup{
	colorlinks=true,
	linkcolor=blue,
	citecolor=blue,
	urlcolor=blue,
}

\title{Risikovergleich gängiger Drogen in Deutschland \\ mit Schwerpunkt Jugendschutz}
\author{Zusammenstellung auf Basis aktueller Forschung}
\date{\today}

\begin{document}
	
	\maketitle
	
	\begin{abstract}
		Die Frage nach der Gefährlichkeit von Drogen ist komplex, da verschiedene Substanzen auf unterschiedliche Weise Schaden anrichten können. Dieses Dokument vergleicht zunächst die drei in Deutschland legal verfügbaren psychoaktiven Substanzen – Tabak, Alkohol und Cannabis – und geht anschließend detailliert auf die Risiken illegaler Drogen wie Opioide, Kokain, Amphetamine und MDMA ein. Es werden wissenschaftliche Erkenntnisse zu akuter Toxizität, Langzeitschäden, Abhängigkeitspotenzial und gesellschaftlichen Folgen zusammengefasst. Ein weiterer Abschnitt beleuchtet die medizinische Anwendung dieser Substanzen. Ergänzt wird der Vergleich durch die zentralen Regelungen des deutschen Jugendschutzes sowie durch einen Überblick über Präventions- und Behandlungsmöglichkeiten.
	\end{abstract}
	
	\section{Einleitung}
	
	In der öffentlichen Diskussion werden Drogen häufig pauschal bewertet, ohne die großen Unterschiede in ihren Risikoprofilen zu berücksichtigen. Während Alkohol und Tabak als legale Genussmittel tief in der Gesellschaft verwurzelt sind, gilt Cannabis oft als relativ harmlos. Ein direkter Vergleich zeigt jedoch, dass Alkohol und Tabak in vielen Kategorien ein höheres Schadenspotenzial aufweisen als Cannabis. Darüber hinaus existieren illegale Substanzen mit teils extrem hohen Risiken, insbesondere für akute Überdosierungen und schwere Abhängigkeiten. Ziel dieses Dokuments ist es, eine differenzierte, wissenschaftlich fundierte Übersicht über die Risiken der verbreitetsten Drogen in Deutschland zu geben – geordnet nach ihrem Rechtsstatus. Ein besonderer Fokus liegt auf dem Schutz von Kindern und Jugendlichen durch gesetzliche Regelungen sowie auf Präventions- und Behandlungsangeboten. Zudem wird das therapeutische Potenzial vieler dieser Substanzen dargestellt, das in kontrollierten medizinischen Kontexten genutzt wird. Bevor die einzelnen Substanzen betrachtet werden, lohnt ein Blick auf die historische Entwicklung des Drogenkonsums und die Entstehung der heutigen Verbotspolitik.
	
	\section{Historische Entwicklung und Gründe für die Prohibition}
	
	Die Frage, warum einige psychoaktive Substanzen legal sind (Alkohol, Tabak, seit 2024 eingeschränkt Cannabis) und andere strikt verboten (Heroin, Kokain, MDMA), lässt sich nur historisch verstehen. Die heutige Rechtslage ist weniger das Ergebnis einer rein medizinischen Risikobewertung als vielmehr das Produkt von kulturellen, wirtschaftlichen und politischen Entwicklungen der letzten 150 Jahre.
	
	\subsection{Jahrtausendealter Drogenkonsum}
	
	Die Nutzung psychoaktiver Substanzen ist keineswegs eine Erscheinung der Moderne. Archäologische Funde belegen, dass Alkohol bereits vor über 9.000 Jahren in China hergestellt wurde, Opium im alten Mesopotamien (um 3.400 v. Chr.) kultiviert wurde und Cannabis sowohl in Zentralasien als auch in Indien seit Jahrtausenden als Rauschmittel und Medizin Verwendung fand. Auch Kokablätter wurden in den Anden seit mindestens 5.000 Jahren gekaut, um Müdigkeit zu bekämpfen. In traditionellen Gesellschaften war der Konsum oft rituell eingebunden und unterlag sozialen Regeln, aber nicht staatlicher Verbote.
	
	\subsection{Vom 19. Jahrhundert zur ersten internationalen Regulierung}
	
	Das 19. Jahrhundert brachte zwei entscheidende Veränderungen: die Isolierung reiner Wirkstoffe (Morphin aus Opium 1804, Kokain 1859, Heroin 1874) und die Einführung des Hypodermikums (Spritze) um 1850. Auf einmal standen hochpotente, injizierbare Substanzen zur Verfügung, was zu einer Zunahme von Abhängigkeit und Überdosierungen führte. Gleichzeitig expandierte der Kolonialhandel: Großbritannien zwang China im Rahmen der Opiumkriege (1839–1842, 1856–1860) zur Öffnung seiner Märkte für britisches Opium aus Indien – ein dunkles Kapitel, das Millionen Chinesen in die Sucht trieb.
	
	Die erste internationale Drogenkonferenz fand 1909 in Shanghai statt, initiiert von den USA, die zunehmend Probleme mit chinesischen Einwanderern und Opiumkonsum sahen. Es folgte die Internationale Opiumkonvention von 1912 im Haag – der erste völkerrechtliche Vertrag, der den nicht-medizinischen Handel mit Opium und Kokain beschränkte.
	
	\subsection{Die US-amerikanische Prohibition als Vorbild}
	
	Die heutige Drogenprohibition hat ihre schärfste Ausprägung in den USA. Der \textit{Harrison Narcotics Tax Act} von 1914 unterwarf Opium und Kokain einer starken Steuer und verschriebene Apotheken als einzige legale Bezugsquelle. Mediziner, die Abhängige behandelten, wurden kriminalisiert. Die Alkoholprohibition (1919–1933) scheiterte zwar spektakulär, schuf aber einen Präzedenzfall für bundesweite Verbote. 
	
	Entscheidend für die Kriminalisierung von Cannabis war die Kampagne von Harry Anslinger, dem ersten Direktor des Federal Bureau of Narcotics (1930–1962). Anslinger verband Cannabis mit rassistischen Stereotypen gegen mexikanische Einwanderer und Jazzmusiker und prägte den Begriff "Marihuana", um die Substanz als fremdartig zu brandmarken. Der \textit{Marihuana Tax Act} von 1937 verbot Cannabis faktisch – nicht aufgrund medizinischer Erkenntnisse, sondern aus moralischer Panik und rassistischen Motiven.
	
	\subsection{Die Ausbreitung des Verbotsmodells auf die Welt}
	
	Nach dem Zweiten Weltkrieg setzten die USA durch ihre Vormachtstellung in den Vereinten Nationen durch, dass das amerikanische Modell global wurde. Das \textit{Einheitsabkommen über die Betäubungsmittel} von 1961 vereinheitlichte die Kontrolle von rund 120 Substanzen (Opium, Kokablätter, Cannabis). Das \textit{Abkommen über psychotrope Substanzen} von 1971 erweiterte die Liste um synthetische Drogen wie Amphetamine, Barbiturate und später MDMA. Die \textit{VN-Konvention gegen den unerlaubten Verkehr mit Suchtstoffen} von 1988 kriminalisierte schließlich nicht nur Besitz und Handel, sondern auch den persönlichen Konsum.
	
	Deutschland setzte diese Vorgaben über das \textit{Betäubungsmittelgesetz (BtMG)} von 1971 (mehrfach novelliert) um, das bis heute die Grundlage der Drogenkriminalisierung bildet – mit Ausnahme der jüngsten Teil-Legalisierung von Cannabis durch das Konsumcannabisgesetz (KCanG) 2024, die einen vorsichtigen Bruch mit der Prohibition darstellt.
	
	\subsection{Gründe für Verbote – eine multidimensionale Betrachtung}
	
	Die historisch gewachsenen Verbote lassen sich aus verschiedenen Blickwinkeln begründen:
	
	\begin{itemize}[leftmargin=*, noitemsep]
		\item \textbf{Gesundheitsschutz:} Viele Drogen (Opioide, Kokain, Methamphetamin) bergen ein hohes Risiko für akute Überdosierungen, schwere körperliche und psychische Langzeitschäden sowie ein starkes Abhängigkeitspotenzial. Verbote sollen den Zugang insbesondere für Kinder und Jugendliche erschweren und die Zahl der Konsumenten reduzieren.
		\item \textbf{Jugendschutz:} Da das Gehirn bis etwa 25 Jahre reift, sind Jugendliche besonders anfällig für suchtbedingte Entwicklungsschäden. Strafrechtliche Verbote flankieren hier die zivilrechtlichen Jugendschutzgesetze.
		\item \textbf{Bekämpfung von Beschaffungskriminalität:} Heroinabhängige finanzieren ihre Sucht oft durch Eigentumsdelikte. Das Verbot von Erwerb und Besitz soll den Markt austrocknen – mit begrenztem Erfolg, wie die Existenz einer florierenden Schwarzmarkthandel zeigt.
		\item \textbf{Internationale Verpflichtungen:} Deutschland ist als UN-Mitgliedstaat an die drei Drogenkonventionen gebunden. Eine einseitige Legalisierung würde völkerrechtliche Verträge verletzen (was die Cannabis-Teil-Legalisierung 2024 bereits de facto tut – sie bewegt sich in einer Grauzone).
		\item \textbf{Kritik an der Prohibition:} Viele Wissenschaftler und Suchtexperten weisen auf die negativen Folgen der Verbote hin: Kriminalisierung von Konsumenten, Verdrängung des Marktes ins organisierte Verbrechen, fehlende Qualitätskontrolle (Verunreinigungen, unterschiedliche Wirkstoffgehalte) und die Stigmatisierung von Abhängigen, die Hilfe scheuen. Diese Kritik hat in der jüngeren Cannabis-Legalisierung in Deutschland, Kanada und Uruguay sowie in der faktischen Entkriminalisierung in Portugal (seit 2001) einen Niederschlag gefunden.
	\end{itemize}
	
	Zusammenfassend lässt sich festhalten: Die heutige Einteilung in „legal“ und „illegal“ ist nicht primär durch medizinische Risikobewertungen bestimmt, sondern durch historische Zufälle, wirtschaftliche Interessen (z.B. Alkohol- und Tabaklobby) und politische Machtverhältnisse. Das Risikoprofil einer Substanz ist jedoch unabhängig von ihrem Rechtsstatus – und genau dieses Risiko wird in den folgenden Abschnitten objektiv dargestellt.
	
	\section{Der „Krieg gegen die Drogen“ – eine kritische Bilanz nach über 50 Jahren}
	
	
	Seit US-Präsident Richard Nixon 1971 den Drogenmissbrauch offiziell zum „Staatsfeind Nr. 1“ erklärte und den „War on Drugs“ ausrief, haben die USA und viele andere Nationen, darunter auch Deutschland, enorme Ressourcen in die repressive Bekämpfung des illegalen Drogenmarktes investiert. Die selbstgesteckten Ziele waren ehrgeizig: Reduzierung von Angebot und Nachfrage, Eindämmung des Schwarzmarktes, Schutz der öffentlichen Gesundheit. Nach über 50 Jahren ist die Bilanz dieser Politik jedoch ernüchternd – und die wissenschaftliche Bewertung fällt nahezu einhellig negativ aus.
	
	\subsection{Kosten: Eine Milliarden-Dollar-Investition ohne Rendite}
	
	Die finanziellen Aufwendungen für den „Krieg gegen die Drogen“ sind immens. In den USA stiegen die Ausgaben von 100 Millionen Dollar im Jahr 1971 auf rund 15 Milliarden Dollar im Jahr 2013, insgesamt über 40 Jahre hinweg auf schätzungsweise eine Billion Dollar.\cite{prohibition-costs} Auch in Deutschland sind die Summen beträchtlich: Pro Jahr werden hierzulande etwa 4 Milliarden Dollar für die Prohibition ausgegeben.\cite{prohibition-costs} Allein der Strafvollzug im Zusammenhang mit Drogendelikten verschlang 2006 rund 850 Millionen Euro.\cite{drug-war-germany-costs} Hinzu kommen Kosten für Polizei, Staatsanwaltschaften, Gerichtsprozesse, Zeugenschutzprogramme und aufwendige Observationen. Legt man eine konservative Ausgabenschätzung für Deutschland über den Zeitraum seit Anfang der 1970er Jahre zugrunde, summieren sich die Kosten auf rund 230 bis 270 Milliarden Euro.\cite{germany-investment}
	
	Diesen enormen Ausgaben stehen jedoch kaum messbare Erfolge gegenüber. Die Nachfrage nach illegalen Drogen wurde nicht gesenkt, die Verfügbarkeit nicht eingeschränkt.\cite{germany-investment} Der illegale Drogenmarkt in Deutschland gehört heute zu den größten in Europa, mit einem jährlichen Umsatz im zweistelligen Milliardenbereich – über die Jahre ist dieser Markt nicht kleiner, sondern größer geworden.\cite{germany-investment} Die Zahl der registrierten Rauschgiftdelikte in der Bundesrepublik hat sich in den letzten 28 Jahren trotz laufend verschärfter Strafgesetze etwa vervierfacht.\cite{nobis-failure}
	
	\subsection{Die unbeabsichtigten Folgen: Schwarzmarkt, Gewalt und Korruption}
	
	Statt den Drogenmarkt auszutrocknen, hat die Prohibition einen riesigen globalen Schwarzmarkt geschaffen.\cite{black-market} Die künstliche Verknappung durch Verbote treibt die Preise in die Höhe – und damit die Gewinnmargen für kriminelle Organisationen. Die Ökonomen Hartwig und Pies kamen bereits 1997 zu dem vernichtenden Urteil: „Prohibition ist nicht nur unnütz, sie schadet sogar.“\cite{spiegel-prohibition} Die hohen Gewinne animierten die Anbieter, mehr Stoff auf den Markt zu drücken, wovon der kontinuierliche Verfall der Heroin-Preise auf das bisher niedrigste Niveau zeugt.\cite{spiegel-prohibition}
	
	Die exorbitanten Gewinnspannen erklären, warum der Konkurrenzkampf verschiedener Drogenanbieter vorzugsweise mit gewalttätigen Mitteln ausgetragen wird.\cite{mexico-drug-war} Die Differenz zwischen dem Endverbraucherpreis und den tatsächlichen Produktionskosten ist bei keiner Güterkategorie größer als bei illegalen Drogen.\cite{mexico-drug-war} Während kolumbianische Bauern 2008 etwa 450 Euro pro Kilo Kokain erhielten, kostete das Kilo in Deutschland schließlich zwischen 25.000 und 50.000 Euro – eine Verhundertfachung des Ertrags.\cite{drug-war-failure}
	
	Die Folgen dieser Ökonomie sind verheerend: In Mexiko erklärte Präsident Calderón 2006 den Drogenkartellen den „Krieg“ und setzte rund 45.000 Soldaten ein – doch der Einsatz des Militärs hat die Lage eher verschlimmert als verbessert.\cite{mexico-drug-war} In Guatemala sterben pro 100.000 Einwohner zwölf Menschen an den Folgen des Drogenkriegs, viermal so viele wie in Mexiko.\cite{drug-war-failure} In autoritären Staaten wie den Philippinen wurden unter Präsident Duterte zwischen 2016 und 2022 über 6.000 Menschen von der Polizei in Anti-Drogen-Einsätzen getötet – ein blutiges Kapitel, das schließlich zur Festnahme Dutertes durch den Internationalen Strafgerichtshof wegen Verbrechen gegen die Menschlichkeit führte.\cite{philippines-drug-war}
	
	\subsection{Gesundheitliche Schäden durch Kriminalisierung}
	
	Ein besonders tragisches Paradox des „Kriegs gegen die Drogen“ besteht darin, dass die repressive Politik selbst zu erheblichen gesundheitlichen Schäden beiträgt. Die auf Konsumenten ausgerichtete Repression behindert Maßnahmen zum Schutz der öffentlichen Gesundheit, die darauf abzielen, HIV/Aids, tödliche Überdosen und andere schädliche Folgen des Drogenkonsums einzudämmen.\cite{nobis-failure} Infolge der Drogenprohibition konnten sich durch die gemeinsame Nutzung verunreinigter Spritzen unter Konsumierenden Krankheiten wie HIV/Aids sowie Hepatitis C ausbreiten.\cite{bpb-drug-control}
	
	Fehlende Qualitätskontrollen auf dem Schwarzmarkt führen dazu, dass Substanzen mit gefährlichen Streckmitteln verunreinigt sind und die tatsächliche Wirkstoffkonzentration stark schwankt – ein erhebliches Risiko für ungewollte Überdosierungen. Die Global Commission on Drug Policy kommt daher zu dem Schluss: Die auf die Drogenkonsumierenden ausgerichtete Repression behindert evidenzbasierte Investitionen in die Verringerung der Nachfrage und die Schadensminderung.\cite{nobis-failure}
	
	\subsection{Kriminalisierung und Stigmatisierung von Konsumenten}
	
	Die Prohibition kriminalisiert in erster Linie die Konsumenten, nicht die Probleme. In Deutschland lag der Anteil der wegen Drogendelikten Inhaftierten 2009 bei 15 Prozent (ohne Beschaffungskriminalität).\cite{drug-war-failure} Die Kriminalisierung erschwert den Zugang zu Hilfsangeboten: Konsumenten scheuen sich aus Angst vor Strafverfolgung, ärztliche oder psychosoziale Hilfe in Anspruch zu nehmen. Auch die Präventionsarbeit wird durch die Illegalität von Drogen behindert.
	
	Der Bochumer Ökonom Karl-Hans Hartwig identifizierte einen weiteren fatalen Mechanismus: Der hohe Endverbraucherpreis versetzt Konsumenten unter einen enormen Finanzierungsdruck, sodass etwa ein Drittel ihres Bedarfs durch Dealen gedeckt wird. Neulinge werden in der Regel von süchtigen Kleinhändlern aus ihrem Bekanntenkreis verführt, der erste Schuss ist meist ein Geschenk. Damit entfällt der erklärte Sinn des Verbots – dass der hohe Preis eine Einstiegshürde für Erstkonsumenten darstellen soll – vollständig. Die Prohibition erreicht das genaue Gegenteil von dem, was eigentlich erwünscht wäre.\cite{spiegel-prohibition}
	
	\subsection{Alternativen: Der portugiesische Weg}
	
	Dass es auch anders geht, zeigt das Beispiel Portugal. Das Land, das Anfang der 2000er Jahre von einer verheerenden Heroinkrise geplagt war – schätzungsweise ein Prozent der Bevölkerung war heroinabhängig, Portugal hatte die höchste HIV-Infektionsrate in der gesamten Europäischen Union – unternahm 2001 einen radikalen Schritt: Es entkriminalisierte den Besitz und Konsum aller illegalen Drogen.\cite{portugal-model}
	
	Wer mit Eigenverbrauchsmengen erwischt wird, erhält keine Strafe mehr, sondern muss vor einer Kommission aus Juristen, Sozialarbeitern und Psychologen erscheinen, die über Behandlung, Schadensbegrenzung und Hilfsdienste informiert.\cite{portugal-model} Drogen sind in Portugal bis heute nicht legal, aber die Entkriminalisierung hat die Drogenpolitik von der Strafjustiz ins Gesundheits- und Sozialwesen verlagert.
	
	Die Ergebnisse sind beeindruckend: Seit der Entkriminalisierung ist der Drogenkonsum allgemein und besonders bei jungen Menschen stark gesunken. Die Zahl der drogenbedingten Todesfälle, Inhaftierungen und HIV-Infektionen ging drastisch zurück.\cite{portugal-results} Portugal hat mit seinem mutigen Experiment gegenüber den Alarmisten im In- und Ausland Recht behalten – der befürchtete massive Anstieg des Drogenkonsums und eine Flut von Drogentouristen blieben aus.\cite{portugal-results} Portugal zeigt, dass eine entkriminalisierende, auf Prävention und Gesundheitsfürsorge ausgerichtete Drogenpolitik nicht humaner, sondern auch effektiver sein kann als die repressive Prohibition.
	
	\subsection{Ausblick: Eine Neubewertung ist überfällig}
	
	Die Bilanz nach über 50 Jahren „Krieg gegen die Drogen“ ist eindeutig: Der Krieg ist gescheitert, schädlich und teuer.\cite{uni-muenster-prohibition} Die selbstgesteckten Ziele wurden verfehlt – stattdessen hat die Prohibition einen blühenden Schwarzmarkt, organisierte Kriminalität, massive Gewalt, Korruption und zusätzliche gesundheitliche Schäden hervorgebracht. Weltweit wächst die Einsicht, dass grundlegende Reformen der Drogenpolitik dringend erforderlich sind.\cite{nobis-failure} Die Teil-Legalisierung von Cannabis in Deutschland, die Entkriminalisierung in Portugal und die Legalisierung in kanadischen und US-amerikanischen Bundesstaaten sind erste Schritte in eine andere Richtung. Doch der „Krieg gegen die Drogen“ ist damit keineswegs beendet – und die kritische Auseinandersetzung mit seinen Folgen bleibt eine zentrale Aufgabe von Wissenschaft und Politik.
	
	\section{Sport und Bildung als zentrale Säulen der Suchtprävention}
	
	Nach der kritischen Betrachtung der repressiven Ansätze rückt nun die Frage in den Fokus, wie Suchtprävention langfristig und nachhaltig gelingen kann. Jenseits von Verboten und Strafverfolgung sind es vor allem die Stärkung von Lebenskompetenzen und die Schaffung gesunder Lebenswelten, die als wirksame Schutzfaktoren vor Suchtentwicklung gelten. Sport und Bildung bieten hierfür ideale Settings, da sie einen großen Teil der Lebenszeit von Kindern und Jugendlichen strukturieren und in ihnen Werte wie Teamgeist, Verantwortung und Selbstwirksamkeit vermitteln können.
	
	\subsection{Bildung: Lebenskompetenzen als Schlüssel zur Prävention}
	
	Die Schule ist als zentraler Lernort und sozialer Lebensraum prädestiniert für systematische und nachhaltige Präventionsarbeit. Das vorrangige Ziel schulischer Suchtprävention ist es, die Lebenskompetenzen von Kindern und Jugendlichen zu stärken\cite{school-prevention}. Aktive und selbstbestimmte Jugendliche, die selbstbewusst sind und Gefühle zulassen können, die angesichts von Enttäuschungen nicht gleich resignieren, sind am wenigsten suchtgefährdet\cite{school-prevention}.
	
	Im Mittelpunkt dieses Ansatzes steht die Förderung von Lebenskompetenzen – also jener Fähigkeiten, die Menschen benötigen, um mit den Herausforderungen des täglichen Lebens erfolgreich umzugehen\cite{life-skills-bzga}. Besonders nachhaltig ist diese Förderung, wenn sie zum Herzstück eines umfassenden Präventionskonzepts in Bildungseinrichtungen wird\cite{life-skills-bzga}. Konkret umfasst dies:
	
	\begin{itemize}[leftmargin=*, noitemsep]
		\item \textbf{Kognitive Dimension:} Aufbau von Faktenwissen über Substanzen und ihre Wirkungen, um sachbezogene Einsichten zu ermöglichen.
		\item \textbf{Wertdimension:} Vermittlung von verbindlichen Werten als Richtlinien des Handelns und die Fähigkeit zur situationsangemessenen Bewertung.
		\item \textbf{Handlungsdimension:} Training von Problemlösefähigkeiten und sozialen Kompetenzen, um in Drucksituationen widerstehen zu können.
	\end{itemize}
	
	Evaluationsergebnisse bestätigen die positive Wirkung von Lebenskompetenzprogrammen in der Suchtprävention\cite{life-skills-evaluation}. Besonders bedeutsam sind der verhinderte Substanzkonsum von Tabak und Alkohol bei jüngeren Schülerinnen und Schülern sowie der verzögerte Konsumbeginn um ein bis zwei Jahre\cite{life-skills-evaluation}. Programme wie „Klasse2000“ oder „HaLT – Hart am Limit“ setzen genau hier an und integrieren die Förderung von Lebenskompetenzen explizit in suchtpräventive Projekte.
	
	\subsection{Sport: Vom Risikofaktor zum Schutzfaktor}
	
	Sportvereine erreichen einen großen Teil der Kinder und Jugendlichen in Deutschland und können unter entsprechenden Bedingungen ein sehr gut geeignetes Setting für eine differenzierte Suchtprävention darstellen\cite{dosb-sport-prevention}. Der Sport vermittelt Erfolgserlebnisse, fördert soziale Kompetenz und Bewältigungsstrategien für Belastungssituationen. Dies fördert den Aufbau eines positiven Selbstkonzepts und Selbstwertgefühls, was Heranwachsenden hilft, Verlockungen und Gruppendruck zu widerstehen\cite{dosb-sport-prevention}.
	
	Allerdings ist die Wirkung von Sport nicht automatisch suchtpräventiv. Studien aus der Schweiz und Deutschland zeigen, dass jugendliche Vereinssportler im Durchschnitt mehr Alkohol konsumieren als Nicht-Sportler – insbesondere nach dem Training oder Wettkampf ist Bier ein verbreitetes Ritual, bei dem jüngere Spieler die Älteren nachahmen\cite{sport-prevention-risks}. Bei Tabak hingegen ist ein umgekehrtes Phänomen erkennbar: Sportler rauchen in der Regel weniger als Nichtsportler, was vor allem für Ausdauersportler gilt\cite{sport-prevention-risks}.
	
	Diese ambivalente Befundlage bedeutet: Sport allein ist keine Suchtvorbeugung\cite{sport-instead-of-addiction}. Erst die gezielte suchtpräventive Arbeit innerhalb der Vereine – etwa durch Verhaltenskodizes, die Förderung von Lebenskompetenzen und die Sensibilisierung der Betreuer – kann das Potenzial des Sports ausschöpfen.
	
	\subsection{Die therapeutische Wirkung von Sport}
	
	Über die primäre Prävention hinaus spielt Sport auch in der Sekundärprävention und Rehabilitation eine wichtige Rolle. In der Suchttherapie kann regelmäßige Bewegung den Behandlungsplan sinnvoll ergänzen und trägt dazu bei, den Erfolg der Suchtbehandlung sowie die körperliche und psychische Gesundheit der Betroffenen zu verbessern\cite{sport-therapy}. Suchtkranke befinden sich in der Regel in einem körperlich desolaten Zustand; Sport hilft ihnen dabei, ein Gefühl für körperliches und psychisches Wohlbefinden zu erlangen\cite{dosb-sport-prevention}.
	
	Nach einer erfolgreichen Behandlung sind Rückfälle keine Seltenheit. Bei Patientinnen und Patienten, die sich aufgrund einer Alkoholkonsumstörung in Behandlung begeben, liegt die Rückfallquote bei bis zu 70 Prozent. Bewegung und Sport können hier als stabilisierender Faktor wirken\cite{sport-therapy}.
	
	\subsection{Integration und Synergien}
	
	Die größte präventive Wirkung entfalten Sport und Bildung nicht isoliert, sondern in Kombination. Die Deutsche Hauptstelle für Suchtfragen (DHS) betont in ihrem aktuellen REITOX-Bericht 2025 die Bedeutung von Evaluation, Vernetzung und dem Transfer bewährter Maßnahmen zur Verbesserung der Wirksamkeit von Suchtprävention\cite{dhs-reitox-2025}. Lebenskompetenzprogramme wirken noch besser in Kombination mit verhältnispräventiven Ansätzen und der gezielten Förderung von Kindern und Jugendlichen mit schlechteren Startbedingungen\cite{life-skills-evaluation}.
	
	Die Förderung der Gesundheitskompetenzen ist in den Bildungsplänen vieler Bundesländer fächerintegrativ verankert – von der ersten Klasse an. Dabei werden sowohl die Vermeidung von gesundheitsschädlichem Verhalten, Bewegungsförderung als auch die Stärkung von psychischer Gesundheit und Resilienz fokussiert\cite{health-literacy-schools}.
	
	\subsection{Praxisbeispiele aus Deutschland}
	
	In der Praxis finden sich zahlreiche Ansätze, die die theoretischen Überlegungen umsetzen:
	
	\begin{itemize}[leftmargin=*, noitemsep]
		\item \textbf{Klasse2000:} Das bundesweit verbreitete Grundschulprogramm fördert Gesundheitsverhalten und Lebenskompetenzen – mit nachgewiesener Wirkung auf den späteren Substanzkonsum.
		\item \textbf{HaLT – Hart am Limit:} Das Programm der Deutschen Hauptstelle für Suchtfragen richtet sich an Jugendliche mit riskantem Alkoholkonsum und setzt auf Frühintervention.
		\item \textbf{„Kinder stark machen“:} Eine Qualifizierungsinitiative der BZgA, die suchtpräventive Arbeit speziell im Kinder- und Jugendsport etabliert.
		\item \textbf{Landesprogramm Bildung und Gesundheit (LaPro BuG) in NRW:} Unterstützt Schulen bei der integrierten Gesundheits- und Qualitätsentwicklung\cite{school-prevention}.
	\end{itemize}
	
	\subsection{Grenzen und Herausforderungen}
	
	Trotz der vielversprechenden Ansätze bleiben Herausforderungen bestehen. Die Umsetzung evidenzbasierter Programme scheitert häufig an begrenzten Ressourcen, fehlender Ausbildung der Multiplikatoren oder mangelnder Verstetigung. Auch die Gefahr einer reinen „Alarmprävention“ – die Jugendlichen lediglich abschrecken will – ist nicht gebannt. Lebenskompetenzprogramme ersetzen keine strukturellen Maßnahmen wie Werbeverbote oder Preisregulierungen; sie können diese aber sinnvoll ergänzen\cite{life-skills-evaluation}.
	
	Sportvereine stehen zudem vor der Herausforderung, ihre eigenen Risiken zu adressieren. Der Sport bietet nicht nur Schutz, sondern kann auch Suchtverhalten fördern – sei es durch die bereits beschriebene Alkoholkultur, durch Doping oder durch die spezifische Gefahr von Snus (Tabakbeutel), das in Sportarten mit skandinavischem Einfluss weit verbreitet ist\cite{sport-prevention-risks}. Die Bekämpfung dieser Doping- und Suchtmentalität ist daher ebenso notwendig wie die Förderung eines gesundheitsbewussten Vereinsklimas.
	
	Zusammenfassend lässt sich festhalten: Sport und Bildung sind keine Patentrezepte, aber unverzichtbare Bausteine einer modernen, evidenzbasierten Suchtprävention. Sie adressieren die Ursachen von Suchtentwicklung – mangelnde Lebenskompetenzen, geringes Selbstwertgefühl, fehlende soziale Einbindung – und nicht nur die Symptome. In Kombination mit verhältnispräventiven Maßnahmen und einer entstigmatisierenden Drogenpolitik können sie einen entscheidenden Beitrag zur Reduzierung substanzbezogener Probleme leisten.
	
	\section{Drogen, Drogenkriminalität und soziale Folgen in Stuttgart und Baden-Württemberg}
	
	Während die bisherigen Abschnitte überwiegend bundesweite und internationale Perspektiven eingenommen haben, rückt dieser Abschnitt die Situation in der Landeshauptstadt Stuttgart und dem Bundesland Baden-Württemberg in den Mittelpunkt. Die Region weist Besonderheiten auf – von der überdurchschnittlichen Belastung durch Rauschgiftkriminalität bis hin zu spezifischen Herausforderungen durch Drogenlieferdienste („Drogentaxis“) und einer vergleichsweise hohen Zahl von Drogentoten.
	
	\subsection{Rauschgiftkriminalität in Baden-Württemberg}
	
	Nach einer Analyse des Landeskriminalamts Baden-Württemberg war die Landeshauptstadt Stuttgart im Jahr 2022 gemessen an den Einwohnerzahlen die deutsche Großstadt mit den höchsten Drogendelikten pro Kopf \cite{stuttgart-drug-delikte}. Dieser Spitzenwert steht unter anderem im Zusammenhang mit neuen Vertriebswegen über das Darknet und Messenger-Dienste wie Telegram oder Instagram, die den Zugang zu Drogen erheblich vereinfachen \cite{stuttgart-drug-delikte}.
	
	Die Polizeiliche Kriminalstatistik Baden-Württembergs für das Jahr 2024 zeigt eine ambivalente Entwicklung: Während die Gesamtkriminalität im Land im fünften Jahr in Folge rückläufig war, bleibt die Rauschgiftkriminalität ein Problembereich. Die Bekämpfung des illegalen Drogenhandels und -konsums hat daher weiterhin eine hohe Priorität \cite{sicherheitsbericht-2025}. Im Jahr 2024 registrierte das Landeskriminalamt Baden-Württemberg einen Anstieg der Handels- und Beschaffungsdelikte im Zusammenhang mit Kokain \cite{stuttgart-drug-delikte}.
	
	\subsection{Drogentote: Stuttgart im Fokus}
	
	Die Landeshauptstadt Stuttgart ist der regionale Hotspot der Drogentoten in Baden-Württemberg. Im Jahr 2025 starben landesweit 191 Menschen an den Folgen ihres Drogenkonsums, davon allein 40 in Stuttgart \cite{drogentote-2025}. Die Stadt hat sich damit im Vergleich zum Vorjahr (2024: 20 Todesfälle in Stuttgart bei 195 landesweit) massiv verschlechtert. Im Jahr 2023 waren es noch 21 Drogentote in Stuttgart bei 141 landesweit \cite{drogentote-2023} – die Stadt hat sich in diesem erschreckenden Trend mittlerweile an die Spitze des Landes gesetzt.
	
	Das Durchschnittsalter der Verstorbenen im Land lag 2025 bei 37,2 Jahren \cite{drogentote-2025}. Ein hoher Anteil der Todesfälle geht auf den gefährlichen Mischkonsum verschiedener Betäubungsmittel oder von Drogen mit Alkohol zurück, der die tödlichen Wirkungen unkalkulierbar verstärken kann \cite{drogentote-2025}. Die häufigste dokumentierte Todesursache bei Einzelsubstanzkonsum war eine Überdosis Kokain \cite{drogentote-2025}.
	
	\subsection{Drogenkonsum unter Jugendlichen und jungen Erwachsenen}
	
	Eine Besorgnis erregende Entwicklung ist der zunehmende Drogenkonsum unter Jugendlichen und jungen Erwachsenen in der Region. Die Drogenaffinitätsstudie 2025 für Baden-Württemberg zeigt einen signifikanten Anstieg insbesondere bei jungen Frauen: Ihr Anteil stieg von 9,7 Prozent auf 18,8 Prozent. Bei jungen Männern gaben 2023 bereits 26,9 Prozent an, in den letzten 30 Tagen eine illegale Droge konsumiert zu haben \cite{drogenaffinitaetsstudie}.
	
	Besonders tragisch ist die Entwicklung bei den Drogentoten: Im Jahr 2025 waren in Baden-Württemberg 15 Verstorbene Heranwachsende und zwei Personen zählten zur Altersgruppe der Jugendlichen \cite{drogentote-2025}. Die zunehmende Verfügbarkeit von Substanzen – nicht zuletzt durch Online-Handel – wird von Suchtexperten als ein wesentlicher Treiber dieser gefährlichen Entwicklung angesehen.
	
	\subsection{Die offene Drogenszene und polizeiliche Maßnahmen in Stuttgart}
	
	Stuttgart verfügt über eine klassische offene Drogenszene. Historisch war der Bereich rund um die Leonhardskirche und das Carl-Zeiss-Planetarium ein bekannter zentraler Treffpunkt \cite{drogenszene-stuttgart}. Die Polizei beobachtet seit längerer Zeit, dass sich die Szene vom Zentrum weg verlagert und auf mehrere Parks verteilt \cite{drogenszene-verlagerung} – ein klassischer Verdrängungseffekt, der die Probleme nicht löst, sondern nur verlagert.
	
	Das Polizeipräsidium Stuttgart setzt zur Bekämpfung der offenen Drogenszene verschiedene Maßnahmen ein, darunter die Einrichtung eines eigenen Drogenkommissariats (KKD), dessen Beamte verdeckt und spezialisiert gegen Dealer vorgehen. Im Rahmen der Schwerpunktbekämpfung wurden im Jahr 2024 mehrfach aufwändige Operationen durchgeführt: Im September 2024 gelang der Polizei ein Schlag gegen die Gewalt- und Drogenkriminalität im Großraum Stuttgart, bei dem sieben Haftbefehle erlassen wurden und Bargeld in Höhe von mehreren Hunderttausend Euro sowie Drogen, darunter Kokain, beschlagnahmt wurden \cite{razzia-stuttgart-sep24}. Im Dezember 2024 wurde ein Drogennetzwerk im Raum Stuttgart zerschlagen, bei dem zehn Personen festgenommen wurden \cite{drogennetzwerk-stuttgart}.
	
	\subsection{Die neue Dimension: Drogenlieferdienste und Darknet}
	
	Eine der größten Herausforderungen für die Stuttgarter Polizei ist die zunehmende Verlagerung des Drogenhandels in digitale Räume. Dealer nutzen inzwischen häufig Darknet-Plattformen oder Messenger-Dienste wie Telegram und Instagram, um neue Konsumenten zu gewinnen \cite{stuttgart-drug-delikte}. Die Bestelldienste werden wegen ihrer schnellen Lieferung (innerhalb weniger Minuten bis einer Stunde) oft als „Drogentaxis“ bezeichnet \cite{stuttgart-drug-delikte}. Ein journalistischer Selbstversuch bestätigte die Erreichbarkeit und Effizienz dieser Drogenlieferdienste in Stuttgart \cite{stuttgart-drug-delikte}.
	
	Die Antwort der Politik auf diese neue Dimension des Handels war ein Landtagsantrag der FDP/DVP-Fraktion vom Juli 2024, der die Landesregierung zu einer umfassenden Bestandsaufnahme aufforderte – unter anderem zu den Vertriebswegen im Darknet und Clearweb, der Bedeutung von Messenger-Diensten, der Verbreitung von Drogentaxis und der Effektivität polizeilicher Ermittlungen \cite{landtag-antrag}. Der Antrag unterstrich die Notwendigkeit, die Ermittlungsstrategien an die neuen Realitäten des digitalen Drogenhandels anzupassen.
	
	\subsection{Drogen in der Justiz und die Folgen für Abhängige}
	
	Die Justiz in Baden-Württemberg ist mit einer Vielzahl von Drogendelikten befasst. Allein die Zahl der wegen Verstößen gegen das Betäubungsmittelgesetz (BtMG) inhaftierten Personen ist signifikant, wobei es sich häufig nicht um Schwerstkriminelle, sondern um abhängige Konsumenten handelt, deren Sucht sie in die Kriminalität treibt.
	
	In Stuttgart gibt es etwa 1.000 Personen, die im Rahmen einer Substitutionsbehandlung Opioid-Ersatzstoffe von Ärzten verschrieben bekommen \cite{substitution-stuttgart}. Das Angebot der ambulanten Suchthilfe in Stuttgart ist vielschichtig und umfasst Prävention, Beratung, Behandlung, Rehabilitation, Schadensminimierung und niedrigschwellige Überlebenshilfen \cite{suchthilfe-stuttgart}. Die Suchtberatung des Klinikums Stuttgart bietet regelmäßige offene Sprechstunden für Bürgerinnen und Bürger an \cite{suchtberatung-stuttgart} – ein wichtiges Element, um den Kontakt zu den Hilfesuchenden herzustellen.
	
	\subsection{Rekordfunde und internationale Dimension}
	
	Die Stuttgarter Polizei und das Landeskriminalamt Baden-Württemberg sind nicht nur mit lokaler, sondern auch mit internationaler Drogenkriminalität befasst. Das Landeskriminalamt war maßgeblich an der Sicherstellung der bislang größten Kokainmenge in der Bundesgeschichte beteiligt: Insgesamt wurden mehr als 35 Tonnen Kokain mit einem Straßenverkaufswert von mehreren Milliarden Euro beschlagnahmt, zwischengelagert und schließlich vernichtet – die Pakete wurden zeitweise im LKA in Bad Cannstatt unter höchster Sicherheitsstufe gelagert \cite{kokainfund-lka}.
	
	Diese Dimension des Handels zeigt, dass Stuttgart und Baden-Württemberg nicht nur ein regionaler Umschlagplatz, sondern auch ein Drehkreuz für internationale Drogenkartelle sind. Die Beschlagnahmungen sind zwar beeindruckende polizeiliche Erfolge, sie stellen aber nur einen Bruchteil des tatsächlich gehandelten Volumens dar und belegen zugleich die andauernde Flut illegaler Substanzen auf dem Schwarzmarkt.
	
	\subsection{Drogen, Drogenkriminalität und soziale Folgen im Landkreis Esslingen}
	
	Während Stuttgart als Großstadt einen eigenen Hotspot darstellt, zeigt der Blick in den Landkreis Esslingen, dass die Probleme der Drogenkriminalität längst auch die Fläche erreicht haben. Gemessen an der Einwohnerzahl ist der Landkreis mit zuletzt 23.074 Straftaten im Jahr 2025 einer der kriminalitätsbelasteten Kreise in der Region. Die Rauschgiftkriminalität hat sich dabei zu einem wachsenden Problemfeld entwickelt.
	
	\subsubsection{Drogentote: Ein Jahr mit sieben Todesfällen}
	
	Im Jahr 2023 waren im Landkreis Esslingen sieben Todesfälle auf den Konsum illegaler Drogen zurückzuführen. Damit belegt der Kreis – gemeinsam mit dem Ortenaukreis – nach der Landeshauptstadt Stuttgart mit 21 Todesfällen den zweiten Platz in der traurigen Statistik Baden-Württembergs \cite{esslingen-drogentote}. Besonders beunruhigend ist laut Fachleuten ein deutlicher Anstieg der Todesfälle durch Kokain und Crack sowie durch den Missbrauch von Substitutionsmitteln aus der ärztlich verordneten Drogenersatztherapie. Auch die Zahl der Todesfälle durch Mischkonsum verschiedener Substanzen wie Heroin, Kokain und Benzodiazepine stieg zuletzt deutlich an.
	
	\subsubsection{Einbrüche, Bandenkriminalität und illegale Plantagen}
	
	Die Polizei im Kreis Esslingen kämpft nicht nur gegen den Kleindeal, sondern auch gegen Bandenkriminalität. Im Januar 2025 nahm die Polizei Esslingen sechs Männer im Alter von 17 bis 26 Jahren fest, denen unter anderem Raub- und Drogendelikte sowie Geiselnahmen vorgeworfen werden. Die Bande soll in wechselnden Besetzungen in Esslingen, Altbach und Stuttgart schwere Raubstraftaten begangen haben: Sie lockte Rauschgifthändler an entlegene Orte, nahm ihnen teils mit vorgehaltener Schusswaffe die Drogen ab, um das Rauschgift dann selbst zu verkaufen. Die Durchsuchungen förderten Bargeld im Wert von mehreren zehntausend Euro, rund 1.000 Tabletten Ecstasy, Schreckschusswaffen sowie scharfe Schusswaffen samt Munition zutage \cite{esslingen-bande}.
	
	Bereits im August 2024 führte die Polizei unter der Federführung der Kriminalpolizei Esslingen einen groß angelegten Einsatz gegen eine Gruppe aus Kirchheim, die bandenmäßig mit Kokain handeln soll. Bei den Durchsuchungen in mehreren Orten des Landkreises, schwerpunktmäßig in Kirchheim unter Teck, stellten die Beamten rund 60 Gramm Kokain, 20 Gramm Haschisch, 60 Gramm Marihuana, 15 Gramm Amphetamin, mehr als 28.000 Euro Bargeld und drei Schreckschusswaffen sicher \cite{esslingen-kokain}.
	
	Auch illegale Cannabis-Plantagen sind im Kreis keine Seltenheit: Im Oktober 2024 hoben Spezialkräfte der Kriminalpolizei eine Plantage in Gewächshäusern im Esslinger Stadtgebiet aus und beschlagnahmten 190 Cannabispflanzen in voller Blüte. In der Wohnung des 37-jährigen Betreibers fanden die Ermittler zudem über zwei Kilogramm Cannabis \cite{esslingen-plantage}.
	
	\subsubsection{Unfälle unter Drogeneinfluss: Cannabis im Straßenverkehr}
	
	Im Jahr 2024 gab es im Kreis Esslingen deutlich mehr Unfälle unter Drogeneinfluss als bislang. Die Zahl der registrierten Unfälle stieg auf 40 – ein Anstieg um 67 Prozent gegenüber dem Vorjahr. Erstmals seit 2020 starben dabei auch Menschen. Rund 80 Prozent der ertappten Verkehrsteilnehmer hatten Cannabis konsumiert. Die Polizei schließt einen Zusammenhang mit der Cannabis-Legalisierung nicht aus und weist darauf hin, dass vielen Verkehrsteilnehmern nicht ausreichend bewusst sei, dass für das Fahren unter Cannabiseinfluss strikte Grenzwerte und drakonische Strafen drohen \cite{esslingen-unfaelle}.
	
	\subsubsection{Suchtprävention und Beratungsangebote}
	
	Der Landkreis Esslingen verfügt mit seiner zentralen Beratungsstelle für Sucht und Prävention über ein vielschichtiges Hilfesystem. Die Beratungsstelle an vier Standorten bietet kostenfreie, schweigepflichtige Beratung zu legalen und illegalen Substanzen sowie zu Glücksspielsucht an und vermittelt bei Bedarf in weiterführende Behandlungs- und Rehabilitationsangebote. Spezielle Angebote richten sich an substituierte Opioidabhängige, die hier eine ergänzende psychosoziale Begleitung erhalten können. Zudem werden Kurse wie HaLT (Hart am Limit), FreD (Frühintervention bei Drogenkonsum) und PAJ (Polizeiliche Anschlussarbeit mit Jugendlichen) angeboten, die zum Teil auf gerichtlicher Auflage beruhen \cite{esslingen-praevention}.
	
	Die Suchtprävention des Landkreises registriert einen deutlichen Wandel im Konsumverhalten: Alkohol und Tabak verlieren bei Jugendlichen an Bedeutung – während der Anteil der 14- bis 17-Jährigen mit regelmäßigem Alkoholkonsum von rund 80 Prozent im Jahr 2001 auf unter 50 Prozent gesunken ist, liegt der Anteil jugendlicher Raucher heute bei nur noch rund sieben Prozent. Dennoch warnt die Beauftragte für Suchtprävention, Christiane Heinze, vor einer vorschnellen Entwarnung: Fast 20 Prozent der Erwachsenen im Kreis hätten einen riskanten oder schädlichen Alkoholkonsum. Neue Herausforderungen ergeben sich durch eine zunehmende Kokainnutzung und die wachsende Mediensucht \cite{esslingen-praevention}.
	
	In den Schulen des Landkreises setzt die Prävention auf verschiedene Bausteine: die Ausbildung von Schülermultiplikatoren für Suchtprävention, die polizeiliche Drogenprävention mit Schwerpunkten auf Cannabis und synthetischen Drogen sowie spezielle Elternabende. Eine regionale Anlaufstelle für diese Maßnahmen ist Christiane Heinze in der Suchtprophylaxe des Landratsamts Esslingen (Telefon 0711 3902-41578, E-Mail suchtprophylaxe@lra-es.de) \cite{esslingen-kontakt}.
	
	\subsubsection{Fazit für den Landkreis Esslingen}
	
	Der Landkreis Esslingen zeigt stellvertretend, dass die Drogenproblematik längst nicht auf Großstädte beschränkt ist. Der Kreis verzeichnet eine hohe Belastung durch Rauschgiftdelikte, gewalttätige Bandenkriminalität, einen bemerkenswerten Anstieg von Unfällen unter Drogeneinfluss und mit sieben Drogentoten im Jahr 2023 einen der höchsten Werte im Land. Die Herausforderungen sind vielfältig: von der Bekämpfung des internationalen Kokainhandels über die Zerschlagung lokaler Bandenstrukturen bis hin zur Sensibilisierung von Autofahrern für die Risiken des Fahrens unter Cannabis. Die bestehenden Präventions- und Beratungsangebote sind gut aufgestellt, stehen aber angesichts des veränderten Konsumverhaltens – sinkende Zahlen bei Alkohol und Tabak, steigende bei Kokain und Mediensucht – vor neuen Aufgaben. Ähnlich wie auf Landesebene bleibt festzuhalten: repressive Maßnahmen allein greifen zu kurz; nur ein integrativer Ansatz aus Polizeiarbeit, Prävention und niedrigschwelliger Suchthilfe kann die Probleme im Kreis Esslingen langfristig in den Griff bekommen.
	
	\subsection{Fazit zur regionalen Situation}
	
	Die regionale Analyse zeigt, dass Stuttgart und Baden-Württemberg in besonderem Maße von den Problemen der Drogenkriminalität betroffen sind. Die Stadt hat nicht nur die höchste Drogendelikte-pro-Kopf-Rate unter deutschen Großstädten, sondern auch die höchste Zahl an Drogentoten im Land. Die Digitalisierung des Handels über Drogentaxis und Darknet-Plattformen stellt die Polizei vor neue Herausforderungen, die mit traditionellen Ermittlungsmethoden allein nicht zu bewältigen sind. Der Landkreis Esslingen zeigt, dass die Probleme längst im Umland angekommen sind und dort ähnlich drastische Formen annehmen. Die sozialen Folgen sind gravierend: Abhängige geraten in die Kriminalität, sterben an Überdosen oder sind auf langwierige Substitutionsbehandlungen angewiesen. Die Präventions- und Hilfsangebote in der Region sind zwar vielfältig, aber die steigenden Fallzahlen zeigen, dass sie allein die Probleme nicht lösen können. Die regionale Situation unterstreicht damit die bundesweiten Befunde: Eine repressive Politik allein ist gescheitert – es bedarf neuer Ansätze, die sowohl die Gesundheit der Konsumenten in den Mittelpunkt stellen als auch die Ursachen der Sucht bekämpfen.
	
	\section{Legale Drogen: Tabak, Alkohol und Cannabis}
	
	Diese drei Substanzen sind in Deutschland für Erwachsene legal erhältlich (Cannabis seit 2024 unter bestimmten Bedingungen). Sie unterscheiden sich jedoch erheblich in ihrem Gefährdungspotenzial.
	
	\subsection{Tabak / Nikotin}
	
	Tabak ist die mit Abstand tödlichste legale Droge in Deutschland. \cite{dhs-tabak}
	
	\begin{itemize}[leftmargin=*, noitemsep]
		\item \textbf{Akute Toxizität:} Gering – eine akute Nikotinvergiftung ist durch normales Rauchen kaum möglich.
		\item \textbf{Organschäden:} Schwerwiegend – chronisch obstruktive Lungenerkrankung (COPD), Lungenemphysem, Herz-Kreislauf-Erkrankungen (Herzinfarkt, Schlaganfall), periphere arterielle Verschlusskrankheit.
		\item \textbf{Krebsrisiko:} Extrem hoch – Tabak ist ein Gruppe-1-Karzinogen und verursacht Lungen-, Mund-, Rachen-, Speiseröhren-, Blasen-, Nieren- und Bauchspeicheldrüsenkrebs.
		\item \textbf{Psychische Risiken:} Hohe Abhängigkeit (körperlich und psychisch); Entzugserscheinungen (Reizbarkeit, Schlafstörungen, Konzentrationsprobleme). Keine direkte Psychosegefahr.
		\item \textbf{Abhängigkeit:} Extrem hoch – etwa 80\% der Raucher sind nikotinabhängig. Die Bindung an die Substanz ist mit der von Heroin vergleichbar.
		\item \textbf{Gesellschaftlicher Schaden:} Sehr hoch – jährlich über 127.000 Todesfälle in Deutschland, volkswirtschaftliche Kosten von ca. 80 Milliarden Euro pro Jahr (durch Produktivitätsverluste, Krankenkosten, Frühberentung).
	\end{itemize}
	
	\subsection{Alkohol}
	
	Alkohol ist ein Zellgift und in nahezu jeder Risikokategorie gefährlicher als Cannabis. \cite{who-carcinogen}
	
	\begin{itemize}[leftmargin=*, noitemsep]
		\item \textbf{Akute Toxizität:} Hoch – eine Alkoholvergiftung kann tödlich enden. In Deutschland gibt es jährlich ca. 44.000 alkoholbedingte Todesfälle. \cite{who-deaths}\cite{dhs-deaths}
		\item \textbf{Organschäden:} Schwer – Leberzirrhose, Pankreatitis, alkoholische Kardiomyopathie, Hirnatrophie.
		\item \textbf{Krebsrisiko:} Eindeutig (Gruppe 1) – Alkohol ist kausal mit mindestens sieben Krebsarten verbunden (Mund, Rachen, Speiseröhre, Leber, Kehlkopf, Darm, Brust). \cite{who-carcinogen}
		\item \textbf{Psychische Risiken:} Alkoholbedingte Demenz, Depressionen, Angststörungen. Körperlicher Entzug kann lebensbedrohlich sein (Delirium tremens).
		\item \textbf{Abhängigkeit:} Hoch (körperlich und psychisch) – etwa 1,6 Millionen Menschen in Deutschland sind alkoholabhängig. \cite{dhs-abhaengigkeit}
		\item \textbf{Gesellschaftlicher Schaden:} Sehr hoch – 57 Milliarden Euro volkswirtschaftliche Kosten pro Jahr, 30\% aller tödlichen Verkehrsunfälle, hohe Gewalt- und Kriminalitätsbelastung. \cite{dhs-costs}
	\end{itemize}
	
	\subsection{Cannabis}
	
	Cannabis weist im Vergleich zu Alkohol und Tabak ein deutlich geringeres Risiko für körperliche Schäden und gesellschaftliche Folgen auf, birgt jedoch spezifische psychische Gefahren. \cite{nida-no-overdose}
	
	\begin{itemize}[leftmargin=*, noitemsep]
		\item \textbf{Akute Toxizität:} Sehr gering – eine tödliche Überdosis ist extrem unwahrscheinlich. In Deutschland sind keine entsprechenden Fälle verzeichnet. \cite{bmg-cannabis}
		\item \textbf{Organschäden:} Gering – Hauptrisiko durch das Rauchen (oft mit Tabak gemischt) führt zu Atemwegsreizungen. Hinweise auf ein erhöhtes Herz-Kreislauf-Risiko, aber kein eindeutiger Nachweis für Krebs. \cite{nida-heart}
		\item \textbf{Krebsrisiko:} Kein eindeutiger Nachweis für Cannabis allein; Risiko durch Tabak-Mischkonsum.
		\item \textbf{Psychische Risiken:} Erhöhtes Psychoserisiko – besonders bei Jugendlichen mit entsprechender Veranlagung (bis zu 11-fach erhöht). \cite{cannabis-psychosis} Hochpotentes Cannabis verdoppelt das Risiko für psychotische Symptome zusätzlich. \cite{cannabis-high-potency} Beeinträchtigung der Gehirnentwicklung bei Jugendlichen.
		\item \textbf{Abhängigkeit:} Mittel, vorwiegend psychisch – etwa 9\% aller Konsumenten entwickeln eine Abhängigkeit, bei Jugendlichen bis zu 17\%. \cite{dhs-cannabis-abhaengigkeit}
		\item \textbf{Gesellschaftlicher Schaden:} Deutlich geringer als bei Alkohol oder Tabak.
	\end{itemize}
	
	\subsection{Vergleichstabelle: Legale Drogen}
	
	\begin{table}[h!]
		\centering
		\small
		\begin{tabularx}{\textwidth}{@{} l X X X @{}}
			\toprule
			\textbf{Kategorie} & \textbf{Tabak / Nikotin} & \textbf{Alkohol} & \textbf{Cannabis} \\ \midrule
			Akute Toxizität & Gering & Hoch (tödliche Überdosis möglich) & Sehr gering (keine tödliche Überdosis) \\
			Organschäden & Schwer (Lunge, Herz-Kreislauf) & Schwer (Leber, Herz, Gehirn, Pankreas) & Gering (Hauptrisiko durch Rauchen) \\
			Krebsrisiko & Extrem hoch (Gruppe 1) & Eindeutig krebserregend (Gruppe 1) & Kein eindeutiger Nachweis \\
			Psychische Risiken & Abhängigkeit, Entzug & Demenz, Depression, lebensgefährlicher Entzug & Erhöhtes Psychoserisiko (besonders bei Jugendlichen) \\
			Abhängigkeit & Extrem hoch (körperlich \& psychisch) & Hoch (körperlich \& psychisch) & Mittel (vorwiegend psychisch, 9-17\%) \\
			Gesellschaftsschaden & Sehr hoch (80 Mrd. €/Jahr) & Sehr hoch (57 Mrd. €/Jahr) & Gering im Vergleich \\
			\bottomrule
		\end{tabularx}
	\end{table}
	
	\section{Illegale Drogen: Detailierte Betrachtung}
	
	Die folgenden Substanzen unterliegen in Deutschland dem Betäubungsmittelgesetz (BtMG). Ihr Erwerb, Besitz und Handel sind für alle Altersgruppen strafbar. Die Risikoprofile variieren stark – von extrem suchtgefährdenden Opioiden bis zu eher moderat gefährlichen Substanzen wie MDMA.
	
	\subsection{Opioide (Heroin, Fentanyl, synthetische Opioide)}
	
	Opioide verursachen den Großteil der drogenbedingten Todesfälle in Deutschland (ca. 80\%). \cite{drugcom-opioid}
	
	\begin{itemize}[leftmargin=*, noitemsep]
		\item \textbf{Akute Toxizität:} Extrem hoch – Atemstillstand bereits bei geringer Überdosierung. Fentanyl ist etwa 50-100 Mal potenter als Heroin.
		\item \textbf{Organschäden:} Schäden durch intravenösen Konsum (Hepatitis C, HIV, Abszesse, Endokarditis). Bei sauberer Substanz und Applikation keine direkte Organtoxizität.
		\item \textbf{Krebsrisiko:} Kein direkt karzinogen.
		\item \textbf{Psychische Risiken:} Schwerste Abhängigkeit (körperlich dominant). Entzugssyndrom mit heftigen körperlichen Symptomen (Schmerzen, Durchfall, Erbrechen, Krämpfe) – jedoch nicht lebensbedrohlich.
		\item \textbf{Abhängigkeit:} Extrem hoch – bereits nach kurzem regelmäßigem Konsum tritt eine starke körperliche Abhängigkeit ein.
		\item \textbf{Gesellschaftlicher Schaden:} Hoch – Beschaffungskriminalität, Überdosis-Todesfälle, hohe Kosten für Substitutionsbehandlung und Notfallmedizin.
	\end{itemize}
	
	\subsection{Kokain und Crack}
	
	Kokain (intranasal) und seine rauchbare Form Crack unterscheiden sich in Suchtpotenzial und gesundheitlichen Folgen. \cite{nutt-scale}
	
	\begin{itemize}[leftmargin=*, noitemsep]
		\item \textbf{Akute Toxizität:} Hoch – Herzinfarkt, Schlaganfall, Krampfanfälle, Hyperthermie. Crack führt zu rascherer und intensiverer Wirkung, damit höheres Überdosisrisiko.
		\item \textbf{Organschäden:} Herz-Kreislauf-Schäden (Kardiomyopathie, Aortendissektion), Nasenseptumperforation bei nasalem Konsum.
		\item \textbf{Krebsrisiko:} Kein eindeutig karzinogen.
		\item \textbf{Psychische Risiken:} Paranoide Psychosen, Aggressivität, schwere Depressionen nach Konsum ("Cocaine crash").
		\item \textbf{Abhängigkeit:} Sehr hoch – Crack hat ein extremes psychisches Abhängigkeitspotenzial; körperliche Entzugserscheinungen sind mild, aber psychischer Craving sehr stark.
		\item \textbf{Gesellschaftlicher Schaden:} Hoch – insbesondere Crack ist mit Beschaffungskriminalität und sozialem Verfall assoziiert.
	\end{itemize}
	
	\subsection{Amphetamine (besonders Crystal Meth)}
	
	Methamphetamin ("Crystal Meth") ist in Deutschland insbesondere in grenznahen Regionen zu Tschechien verbreitet und führt zu rascher körperlichem und psychischem Verfall. \cite{dhs-crystal}
	
	\begin{itemize}[leftmargin=*, noitemsep]
		\item \textbf{Akute Toxizität:} Hoch – Hyperthermie (Körpertemperatur >40°C), Kreislaufversagen, Krampfanfälle, Rhabdomyolyse.
		\item \textbf{Organschäden:} Schwere Neurotoxizität (Schädigung von Dopamin- und Serotoninnervonen), schwere Zahnschäden ("Meth-Mund"), Kachexie (starker Gewichtsverlust).
		\item \textbf{Krebsrisiko:} Kein direkt karzinogen.
		\item \textbf{Psychische Risiken:} Amphetamin-Psychose (paranoid, halluzinatorisch), anhaltende kognitive Defizite, starke Stimmungsschwankungen.
		\item \textbf{Abhängigkeit:} Sehr hoch – rasche Toleranzentwicklung führt zu Dosissteigerung; schwerer Entzug mit Erschöpfung, Depression, starkem Craving.
		\item \textbf{Gesellschaftlicher Schaden:} Regional sehr hoch (Sachsen, Bayern, Berlin) – Beschaffungskriminalität, Verwahrlosung, Überlastung des Gesundheits- und Justizsystems.
	\end{itemize}
	
	\subsection{MDMA (Ecstasy)}
	
	MDMA hat ein im Vergleich zu anderen illegalen Drogen geringeres Sucht- und Überdosisrisiko, dennoch bestehen erhebliche Gefahren durch Verunreinigungen und falsches Konsumverhalten. \cite{drugcom-mdma}
	
	\begin{itemize}[leftmargin=*, noitemsep]
		\item \textbf{Akute Toxizität:} Mittel – Hyperthermie, Hyponatriämie (durch zu viel Wasser ohne Elektrolyte), Serotoninsyndrom. Todesfälle sind selten, kommen aber vor (v.a. durch Überhitzung auf Partys).
		\item \textbf{Organschäden:} Mögliche Neurotoxizität bei hochdosiertem oder häufigem Konsum (Langzeitschäden an Serotoninsystemen). Leberbelastung.
		\item \textbf{Krebsrisiko:} Kein Nachweis.
		\item \textbf{Psychische Risiken:} Angst- und Panikattacken, Depersonalisationserscheinungen, "Midweek-Blues" (serotonerger Entzug mit depressiver Verstimmung).
		\item \textbf{Abhängigkeit:} Gering – keine körperliche Abhängigkeit; psychische Abhängigkeit ist selten und meist mit partyspezifischen Kontexten verbunden.
		\item \textbf{Gesellschaftlicher Schaden:} Gering im Vergleich zu Opioiden oder Kokain – Gefahren durch verunreinigte Pillen (Streckung mit PMMA, Amphetaminen, synthetischen Cathinonen).
	\end{itemize}
	
	\subsection{Vergleichstabelle: Illegale Drogen}
	
	\begin{table}[h!]
		\centering
		\small
		\renewcommand{\arraystretch}{1.2}
		\begin{tabularx}{\textwidth}{@{} l X X X X X @{}}
			\toprule
			\textbf{Kategorie} & \textbf{Opioide (Heroin)} & \textbf{Kokain / Crack} & \textbf{Amphetamine (Crystal Meth)} & \textbf{MDMA (Ecstasy)} \\ \midrule
			Akute Toxizität & Extrem hoch (Atemstillstand) & Hoch (Herzinfarkt, Schlaganfall) & Hoch (Hyperthermie, Kreislaufversagen) & Mittel (Hyperthermie, Hyponatriämie) \\
			Organschäden & Schäden durch Injektion & Herz-Kreislauf, Nasenseptum & Neurotoxizität, Zahnschäden & Mögliche Neurotoxizität, Leber \\
			Krebsrisiko & Kein direkt karzinogen & Kein eindeutig karzinogen & Kein direkt karzinogen & Kein Nachweis \\
			Psychische Risiken & Schwerste Abhängigkeit & Psychose, Aggressivität, Depression & Psychose, kognitive Defizite & Angst, Panik, Midweek-Blues \\
			Abhängigkeit & Extrem hoch (körperlich dominant) & Sehr hoch, besonders Crack & Sehr hoch (rasche Toleranz) & Gering (keine körperliche) \\
			Gesellschaftsschaden & Hoch (Kriminalität, Todesfälle) & Hoch (Beschaffungskriminalität) & Regional sehr hoch & Gering im Vergleich \\
			\bottomrule
		\end{tabularx}
	\end{table}
	
	\section{Jugendschutz in Deutschland}
	
	Der Gesetzgeber hat für Kinder und Jugendliche gestufte Schutzbestimmungen erlassen, um die gesundheitlichen Risiken und die Entwicklungsschäden durch psychoaktive Substanzen zu minimieren. Die wichtigsten Regelungen im Überblick:
	
	\subsection{Alkohol (\S~9 JuSchG)}
	\begin{itemize}[leftmargin=*, noitemsep]
		\item \textbf{Unter 14 Jahren:} Abgabe und Konsum jeglicher alkoholischer Getränke verboten.
		\item \textbf{14–15 Jahre:} Erlaubt sind Bier, Wein und Sekt ausschließlich in Begleitung einer personensorgeberechtigten Person.
		\item \textbf{16–17 Jahre:} Bier, Wein und Sekt ohne Begleitung erlaubt; branntweinhaltige Getränke (Spirituosen, alkoholhaltige Mixgetränke) bleiben verboten.
		\item \textbf{Ab 18 Jahren:} Keine Beschränkungen.
	\end{itemize}
	
	\subsection{Tabak und nikotinhaltige Erzeugnisse (\S~10 JuSchG)}
	\begin{itemize}[leftmargin=*, noitemsep]
		\item \textbf{Unter 18 Jahren:} Abgabe und Konsum von Tabakwaren, E-Zigaretten, E-Shishas sowie nikotinhaltigen Liquids grundsätzlich verboten. Das Verbot gilt auch für nikotinfreie E-Zigaretten, sofern diese zum Verdampfen von Liquids bestimmt sind.
	\end{itemize}
	
	\subsection{Cannabis (Konsumcannabisgesetz – KCanG)}
	\begin{itemize}[leftmargin=*, noitemsep]
		\item \textbf{Unter 18 Jahren:} Absolutes Cannabisverbot. Erwerb, Besitz und Anbau sind für Minderjährige strafbar. Die Weitergabe an Minderjährige wird mit Freiheitsstrafe nicht unter zwei Jahren bestraft (\S~34 KCanG).
		\item Der Konsum in unmittelbarer Gegenwart von Minderjährigen sowie in einem Umkreis von 100 Metern zu Schulen, Kindertagesstätten und Spielplätzen ist für Erwachsene verboten (\S~5 KCanG).
		\item Verstöße von Jugendlichen werden an das Jugendamt gemeldet und können zur Anordnung von Frühinterventionsmaßnahmen führen.
	\end{itemize}
	
	\subsection{Illegale Drogen (Betäubungsmittelgesetz – BtMG)}
	\begin{itemize}[leftmargin=*, noitemsep]
		\item Substanzen wie Heroin, Kokain, Amphetamine, Methamphetamin und MDMA unterliegen dem BtMG. Erwerb, Besitz und Handel sind für alle Altersgruppen strafbar.
		\item Die Abgabe von Betäubungsmitteln an Personen unter 18 Jahren oder das Bestimmen eines Minderjährigen zum unerlaubten Umgang mit BtM wird als besonders schwerer Fall geahndet (\S~29a Abs.~1 Nr.~1, \S~30 Abs.~1 Nr.~2 BtMG) und mit Freiheitsstrafe nicht unter einem Jahr bzw. nicht unter zwei Jahren bestraft.
	\end{itemize}
	
	\section{Behandlung und Prävention}
	
	Neben gesetzlichen Verboten und Altersgrenzen sind gezielte Präventionsmaßnahmen und ein niedrigschwelliges Behandlungssystem zentrale Säulen der Suchtpolitik. Sie zielen darauf ab, den Einstieg in den Konsum zu verzögern bzw. zu verhindern und betroffenen Personen einen Ausweg aus der Abhängigkeit zu ermöglichen\cite{dhs-praevention}.
	
	\subsection{Präventionsstrategien}
	\begin{itemize}[leftmargin=*, noitemsep]
		\item \textbf{Primärprävention:} Universelle Aufklärung in Schulen, Kindertagesstätten und der Jugendhilfe (z.B. das Programm „Klasse2000“ oder „HaLT – Hart am Limit“). Dazu gehören auch strukturelle Maßnahmen wie Werbebeschränkungen, Steuerpolitik und die Begrenzung von Verkaufsstellen\cite{dhs-praevention}.
		\item \textbf{Selektive Prävention:} Angebote für Risikogruppen, etwa Kinder suchtkranker Eltern, Jugendliche mit Verhaltensauffälligkeiten oder Auszubildende in risikobehafteten Branchen (Gastronomie, Baugewerbe).
		\item \textbf{Indizierte Prävention:} Frühintervention bei erstem problematischem Konsum, z.B. durch Kurzberatung in der Hausarztpraxis, im Jugendamt oder über die Drogenhilfe. Hierzu zählen auch jugendgerichtsnahe Maßnahmen wie die Teilnahme an suchtpräventiven Kursen nach \S~10 JGG.
	\end{itemize}
	
	\subsection{Behandlungsmöglichkeiten}
	\begin{itemize}[leftmargin=*, noitemsep]
		\item \textbf{Entgiftung (qualifizierter Entzug):} Medizinisch überwachte Ausleitung bei körperlicher Abhängigkeit (Alkohol, Opioide, Benzodiazepine). Notwendig wegen möglicher Komplikationen wie Delirium tremens oder Krampfanfällen. Stationäre Entgiftung in Krankenhäusern oder spezialisierten Abteilungen.
		\item \textbf{Psychosoziale Therapien:} Verhaltenstherapie, motivierende Gesprächsführung, kontingenzmanagementbasierte Ansätze. Sowohl ambulant als auch stationär verfügbar.
		\item \textbf{Medikamentengestützte Behandlung:} 
		\begin{itemize}
			\item \textit{Opioide:} Substitutionsbehandlung mit Methadon, Buprenorphin oder retardiertem Morphin (substitutionsgestützte Behandlung).
			\item \textit{Alkohol:} Rückfallprophylaktika (Naltrexon, Nalmefen, Acamprosat) sowie aversive Substanzen (Disulfiram) unter Aufsicht.
			\item \textit{Tabak:} Nikotinersatztherapie (Pflaster, Kaugummi), Vareniclin, Bupropion.
		\end{itemize}
		\item \textbf{Rehabilitation:} Stationäre oder ambulante Langzeittherapie in spezialisierten Einrichtungen (Fachkliniken für Sucht). Ziel ist die dauerhafte Abstinenz oder die Reduktion von Konsumfolgen.
		\item \textbf{Selbsthilfe und Beratung:} Anonyme Alkoholiker, Kreuzbund, Guttempler, cannabisbezogene Selbsthilfegruppen sowie Suchtberatungsstellen der Caritas, Diakonie und kommunaler Träger.
	\end{itemize}
	
	Eine frühzeitige Intervention verbessert die Prognose erheblich. Die gesetzlichen Krankenkassen übernehmen bei gesicherter Diagnose die Kosten für Entgiftung und Rehabilitation\cite{dhs-behandlung}.
	
	\section{Medizinische Anwendung von Drogen}
	
	Die bisherige Betrachtung fokussierte auf die Risiken psychoaktiver Substanzen. Viele dieser Stoffe besitzen jedoch ein erhebliches therapeutisches Potenzial und werden seit Jahrzehnten oder werden zunehmend in kontrollierten medizinischen Kontexten eingesetzt. Der folgende Abschnitt beleuchtet den Stand der medizinischen Anwendung der zuvor besprochenen Substanzen.
	
	\subsection{Cannabis: Etablierte Medizin bei strengen Regeln}
	
	Medizinisches Cannabis hat inzwischen einen festen Platz in der deutschen Therapielandschaft, insbesondere bei chronischen Schmerzen, Spastiken (z.B. bei Multipler Sklerose), Übelkeit und Appetitlosigkeit (z.B. unter Chemotherapie). Seit 2017 können Ärzte es auf Kassenrezept verschreiben, wenn eine schwerwiegende Erkrankung vorliegt, keine Alternative verfügbar ist und eine Verbesserung zu erwarten ist. Aktuelle Daten zeigen, dass über drei Viertel der Verordnungen auf chronische Schmerzen entfallen. Aufgrund eines starken Anstiegs der Verordnungen, besonders über telemedizinische Plattformen, wurden die Regeln 2025 verschärft: Eine Erstverschreibung ist nun nur noch nach persönlichem Arztkontakt möglich, und der Versandhandel mit Cannabisblüten ist verboten.\cite{medical-cannabis}
	
	\subsection{Ketamin \& Esketamin: Durchbruch in der Depressionsbehandlung}
	
	Ein Meilenstein der letzten Jahre ist die Zulassung von Esketamin (einem Ketamin-Abkömmling) als Nasenspray zur Behandlung schwerer, therapieresistenter Depressionen. Dieser Wirkstoff ist ein Paradebeispiel für den Wandel einer Partydroge zum Therapeutikum. Im Gegensatz zu herkömmlichen Antidepressiva, die oft Wochen brauchen, kann die antidepressive Wirkung von Esketamin innerhalb weniger Stunden eintreten, was insbesondere in akuten suizidalen Krisen lebensrettend sein kann. Die Behandlung findet unter strenger ärztlicher Aufsicht statt und ist Patient:innen vorbehalten, bei denen andere Therapien versagt haben.\cite{esketamine}
	
	\subsection{MDMA, Psilocybin \& LSD: Die Renaissance der Psychedelika}
	
	Eine der aufregendsten Entwicklungen ist die Wiederentdeckung von Psychedelika für die Psychotherapie. \textbf{MDMA} wird in klinischen Studien mit großem Erfolg zur Behandlung von posttraumatischen Belastungsstörungen (PTBS) eingesetzt. Es wirkt angstlösend und fördert die therapeutische Beziehung, wodurch Trauma-Erinnerungen sicherer verarbeitet werden können. \cite{mdma-ptsd} Der Wirkstoff aus „Magic Mushrooms“, \textbf{Psilocybin}, fördert die Bildung neuer Nervenverbindungen (Neuroplastizität). Studien zeigen sein Potenzial bei Depressionen, Angstzuständen, Suchterkrankungen und Essstörungen. \textbf{LSD} wird ähnlich erforscht, unter anderem bei Clusterkopfschmerz und Suchterkrankungen. In Deutschland ist eine Psilocybin-Therapie seit Kurzem in Härtefällen (Compassionate Use) an spezialisierten Zentren möglich.\cite{psychedelics}
	
	\subsection{Amphetamine: Standard bei ADHS und Narkolepsie}
	
	Amphetamine haben eine etablierte Rolle in der Schulmedizin. Sie sind ein wichtiger Bestandteil in der Behandlung der Aufmerksamkeitsdefizit-Hyperaktivitätsstörung (ADHS), insbesondere bei schweren Verläufen, die auf andere Medikamente nicht ansprechen. Zudem werden sie zur Therapie der Narkolepsie (krankhafte Schlafsucht) eingesetzt. Sie wirken bei diesen Erkrankungen paradoxerweise beruhigend und strukturierend.\cite{amphetamine-adhd}
	
	\subsection{Opioide: Unverzichtbare, aber risikobehaftete Schmerzmittel}
	
	Trotz ihrer Gefahren als illegale Drogen sind Opioide aus der modernen Schmerzmedizin nicht wegzudenken. Sie sind bei starken und stärksten Schmerzen, etwa bei Tumorerkrankungen, unverzichtbar. Aufgrund ihres hohen Suchtpotenzials unterliegen ihre Verschreibung und Abgabe in Deutschland strengsten gesetzlichen Regelungen.\cite{opioid-pain}
	
	\subsection{Kokain: Historisches Relikt in der Medizin}
	
	Die medizinische Bedeutung von Kokain ist heute fast vollständig von historischem Wert. Es war das erste brauchbare Lokalanästhetikum, das in den 1880er Jahren eingeführt wurde. Es diente chemisch als Vorbild für viele moderne, sicherere Lokalanästhetika und wird heute nur noch in sehr seltenen Ausnahmefällen (z. B. bei bestimmten Nasen-Operationen) eingesetzt.\cite{cocaine-medical}
	
	\section{Fazit}
	
	Die Gegenüberstellung zeigt, dass die legalen Drogen Tabak und Alkohol in Deutschland die größten gesundheitlichen und gesellschaftlichen Schäden verursachen – weit mehr als Cannabis. Unter den illegalen Drogen stellen Opioide wegen ihrer extrem hohen akuten Toxizität und des starken körperlichen Abhängigkeitspotenzials die größte Gefahr dar, gefolgt von Crack und Methamphetamin. MDMA hat im Vergleich ein geringeres Sucht- und Überdosisrisiko, bleibt jedoch aufgrund von Verunreinigungen und risikoreichem Konsumverhalten (Hyperthermie, Hyponatriämie) problematisch. 
	
	Die kritische Bilanz des „Kriegs gegen die Drogen“ zeigt, dass die repressive Verbotspolitik ihre selbstgesteckten Ziele verfehlt hat: Statt Angebot und Nachfrage zu reduzieren, hat sie einen blühenden Schwarzmarkt, organisierte Kriminalität, massive Gewalt und zusätzliche gesundheitliche Schäden hervorgebracht. Das portugiesische Beispiel belegt, dass eine entkriminalisierende, auf Prävention und Gesundheitsfürsorge ausgerichtete Drogenpolitik nicht nur humaner, sondern auch effektiver sein kann.
	
	Gleichzeitig zeigt der medizinische Blick, dass viele dieser Substanzen ein beachtliches therapeutisches Potenzial entfalten – von Cannabis als etabliertem Schmerzmittel über Ketamin als rapidem Antidepressivum bis hin zu psychedelischen Substanzen in der Psychotherapieforschung. Diese Doppelnatur – hohes Missbrauchs- und Schadenspotenzial bei gleichzeitiger medizinischer Nutzbarkeit – erfordert eine differenzierte Betrachtung und eine regulatorische Balance zwischen strikter Kontrolle und therapeutischer Verfügbarkeit.
	
	Keine der genannten Substanzen ist harmlos. Besonders für Kinder und Jugendliche bergen alle psychoaktiven Substanzen erhebliche Entwicklungsrisiken. Die deutschen Jugendschutzgesetze definieren daher gestaffelte Altersgrenzen für legale Drogen und sehen für illegale Drogen strikte Verbote mit empfindlichen Strafen bei Abgabe an Minderjährige vor. Ergänzend leisten Präventionsprogramme und ein vielfältiges Behandlungssystem einen wichtigen Beitrag zur Reduzierung substanzbezogener Probleme. Der sicherste Konsum ist in jedem Alter der Verzicht.
	
	\newpage
	\begin{thebibliography}{99}
		\bibitem{who-deaths} Weltgesundheitsorganisation (WHO), 2025: \href{https://www.who.int/europe/de/news/item/24-12-2025-1-in-3-injury-deaths-caused-by-alcohol--who-europe-reiterates-the-harms}{WHO/Europa weist erneut auf die schädlichen Folgen hin}.
		\bibitem{dhs-deaths} Deutsche Hauptstelle für Suchtfragen (DHS), \textit{Jahrbuch Sucht 2026}, Kapitel 2.1.
		\bibitem{nida-no-overdose} National Institute on Drug Abuse (NIDA); zitiert nach \href{https://www.dw.com/de/faktencheck-wie-ungef%C3%A4hrlich-ist-cannabis/a-65853237}{Deutsche Welle (DW), Faktencheck: Wie (un)gefährlich ist Cannabis?}
		\bibitem{bmg-cannabis} Bundesministerium für Gesundheit: \href{https://www.bundesgesundheitsministerium.de/themen/cannabis/faq-cannabis/ueberdosierung.html}{Cannabis: Fragen und Antworten zur Überdosierung}.
		\bibitem{who-carcinogen} Weltgesundheitsorganisation (WHO), 2021: \href{https://www.who.int/europe/de/home/20-10-2021-alcohol-is-one-of-the-biggest-risk-factors-for-breast-cancer}{Alkohol ist einer der größten Risikofaktoren für Brustkrebs}; IARC-Klassifizierung von 1988.
		\bibitem{nida-heart} National Institute on Drug Abuse (NIDA): \href{https://nida.nih.gov/publications/research-reports/marijuana/cardiovascular-health-effects}{Cardiovascular Health Effects of Marijuana}.
		\bibitem{cannabis-psychosis} André McDonald et al., 2025: Studie zum 11-fach erhöhten Psychoserisiko bei jugendlichen Cannabiskonsumenten, zitiert nach \href{https://www.drugcom.de/news/jugendliche-die-cannabis-konsumieren-haben-ein-erhoehtes-risiko-fuer-psychose/}{drugcom.de}.
		\bibitem{cannabis-high-potency} Lindsey Hines et al., 2024: Längsschnittstudie zu hochpotentem Cannabis und Psychoserisiko, zitiert nach \href{https://www.drugcom.de/news/hochpotenter-cannabis-verdoppelt-risiko-fuer-psychotische-symptome/}{drugcom.de}.
		\bibitem{dhs-abhaengigkeit} Deutsche Hauptstelle für Suchtfragen (DHS): \href{https://www.dhs.de/sucht/zahlen-und-fakten.html}{Zahlen und Fakten zu Alkoholabhängigkeit}.
		\bibitem{dhs-cannabis-abhaengigkeit} Deutsche Hauptstelle für Suchtfragen (DHS), \textit{Jahrbuch Sucht 2025}, Kapitel 2.5; Techniker Krankenkasse: \href{https://www.tk.de/techniker/gesundheit-und-medizin/erste-hilfe-und-gesundheitsvorsorge/sucht-und-abhaengigkeit/cannabis/cannabis-abhaengigkeitspotenzial-2075314}{Cannabis-Abhängigkeitspotenzial}.
		\bibitem{dhs-costs} Deutsche Hauptstelle für Suchtfragen (DHS): \href{https://www.dhs.de/sucht/zahlen-und-fakten.html}{Volkswirtschaftliche Kosten des Alkoholkonsums}.
		% Quellen für erweiterte Substanzen
		\bibitem{euda-2025} Europäische Beobachtungsstelle für Drogen und Drogensucht (EUDA), \textit{Europäischer Drogenbericht 2025}.
		\bibitem{dhs-drugprofiles} Deutsche Hauptstelle für Suchtfragen (DHS): \href{https://www.dhs.de/suechte}{Übersicht zu Stoffgruppen und Abhängigkeitsformen}.
		\bibitem{dhs-tabak} Deutsche Hauptstelle für Suchtfragen (DHS): \href{https://www.dhs.de/suechte/tabak}{Tabak – Zahlen und Fakten}.
		\bibitem{drugcom-opioid} \href{https://www.drugcom.de/drogen/opioide/}{drugcom.de: Opioide}.
		\bibitem{nutt-scale} David Nutt et al., 2010: Drug harms in the UK: a multicriteria decision analysis, \textit{The Lancet}.
		\bibitem{dhs-crystal} Deutsche Hauptstelle für Suchtfragen (DHS): \href{https://www.dhs.de/suechte/crystal-meth}{Crystal Meth – Basisinformationen}.
		\bibitem{drugcom-mdma} \href{https://www.drugcom.de/drogen/ecstasy/}{drugcom.de: Ecstasy (MDMA)}.
		% Quellen für Behandlung und Prävention
		\bibitem{dhs-praevention} Deutsche Hauptstelle für Suchtfragen (DHS): \href{https://www.dhs.de/praevention}{Prävention von Sucht}.
		\bibitem{dhs-behandlung} Deutsche Hauptstelle für Suchtfragen (DHS): \href{https://www.dhs.de/hilfe-und-therapie}{Hilfe und Therapie bei Sucht}.
		% Neue Quellen für medizinische Anwendung
		\bibitem{medical-cannabis} Bundesinstitut für Arzneimittel und Medizinprodukte (BfArM): \href{https://www.bfarm.de/DE/Bundesopiumstelle/Cannabis/Medizinalcannabis/_node.html}{Medizinalcannabis}.
		\bibitem{esketamine} Arzneimittelkommission der deutschen Ärzteschaft: \href{https://www.akdae.de/Arzneimitteltherapie/NA/A-2021-12/}{Esketamin bei therapieresistenter Depression}.
		\bibitem{mdma-ptsd} Multidisciplinary Association for Psychedelic Studies (MAPS): \href{https://maps.org/mdma/}{MDMA-assisted therapy for PTSD}.
		\bibitem{psychedelics} Deutsche Gesellschaft für Psychiatrie und Psychotherapie, Psychosomatik und Nervenheilkunde (DGPPN): Positionspapier zur Psychedelika-Forschung, 2024.
		\bibitem{amphetamine-adhd} Leitlinie der Deutschen Gesellschaft für Kinder- und Jugendpsychiatrie (DGKJP): ADHS im Kindes- und Jugendalter, 2023.
		\bibitem{opioid-pain} Leitlinie der Deutschen Schmerzgesellschaft: Langzeitanwendung von Opioiden bei nicht-tumorbedingten Schmerzen, 2022.
		\bibitem{cocaine-medical} Bundesopiumstelle: \href{https://www.bfarm.de/DE/Bundesopiumstelle/Betaeubungsmittel/BtM-Monographie/_node.html}{Monographie Kokain (medizinische Verwendung)}.
		% Neue Quellen für den Abschnitt "Krieg gegen die Drogen"
		\bibitem{prohibition-costs} Krimpedia: \href{https://www.krimpedia.de/index.php?title=Kritik_der_Prohibition&oldid=68308}{Kritik der Prohibition}.
		\bibitem{drug-war-germany-costs} ik-armut.de: \href{https://ik-armut.de/archiv/inhalt/Krieg\%20gegen\%20Drogen\%20ist\%20gescheitert.htm}{Krieg gegen Drogen ist gescheitert}.
		\bibitem{germany-investment} Stephan König: \href{https://krautinvest.de/koenigs-kolumne-3-geschlossen-wegen-ineffizienz-war-on-drugs/}{König's Kolumne \#3: Geschlossen wegen Ineffizienz – War on Drugs}.
		\bibitem{nobis-failure} Dr. Frank Nobis: \href{http://strafverteidigervereinigungen.de/freispruch/texte/nobis_h2_warondrugs_prn.html}{Ende der Märchenstunde – Weltweit wächst die Einsicht in das Scheitern des War on Drugs}.
		\bibitem{black-market} D+C: \href{https://dandc.eu/de/article/drogen-und-politik-aktuelle-krise-und-historischer-hintergrund}{„Freihandel“ im 19. Jahrhundert}.
		\bibitem{spiegel-prohibition} Spiegel: \href{https://www.spiegel.de/politik/prohibition-ist-nicht-nur-unnuetz-sie-schadet-a-11ecab92-0002-0001-0000-000008651186}{„Prohibition ist nicht nur unnütz, sie schadet“}.
		\bibitem{mexico-drug-war} bpb: \href{https://www.bpb.de/shop/zeitschriften/apuz/33088/calderons-gescheiterter-feldzug-gegen-die-drogenkartelle/}{Calderons gescheiterter Feldzug gegen die Drogenkartelle}.
		\bibitem{drug-war-failure} ik-armut.de: \href{https://ik-armut.de/archiv/inhalt/Krieg\%20gegen\%20Drogen\%20ist\%20gescheitert.htm}{Krieg gegen Drogen ist gescheitert}.
		\bibitem{philippines-drug-war} bpb: \href{https://www.bpb.de/shop/zeitschriften/apuz/philippinen-2025/571160/gewalt-als-vermeintliche-loesung/}{Gewalt als vermeintliche Lösung – Dutertes Drogenkrieg auf den Philippinen}.
		\bibitem{bpb-drug-control} bpb: \href{https://www.bpb.de/kurz-knapp/taegliche-dosis-politik/573244/im-fokus-drogenkontrolle/}{Im Fokus: Drogenkontrolle}.
		\bibitem{portugal-model} The Portugal News: \href{http://new.theportugalnews.com/de/nachrichten/2025-11-04/krieg-gegen-drogen-portugal-hat-vielleicht-die-antwort/624949}{Krieg gegen Drogen? Portugal hat vielleicht die Antwort}.
		\bibitem{portugal-results} Telepolis: \href{https://www.telepolis.de/features/15-Jahre-entkriminalisierte-Drogenpolitik-in-Portugal-3224495.html}{15 Jahre entkriminalisierte Drogenpolitik in Portugal}.
		\bibitem{uni-muenster-prohibition} Universität Münster: \href{https://www.uni-muenster.de/LearnWeb/learnweb2/course/info.php?id=59912}{Kursinformation: Drogenprohibition und die sozio-ökonomischen Folgen}.
		% Neue Quellen für Sport und Bildung
		\bibitem{school-prevention} Ministerium für Schule und Bildung des Landes Nordrhein-Westfalen: \href{https://www.schulministerium.nrw/suchtpraevention-der-schule}{Suchtprävention in der Schule}.
		\bibitem{life-skills-bzga} Bundeszentrale für gesundheitliche Aufklärung (BZgA): \href{https://leitbegriffe.bioeg.de/alphabetisches-verzeichnis/lebenskompetenzen-und-kompetenzfoerderung/}{Lebenskompetenzen und Kompetenzförderung}, 2024.
		\bibitem{life-skills-evaluation} Fachstelle für Suchtprävention Berlin: \href{https://www.berlin-suchtpraevention.de/corona-krise-3-suchtpraevention-lebenskompetenz/}{Suchtprävention fördert Lebenskompetenzen – explizit}, 2020.
		\bibitem{dosb-sport-prevention} Deutscher Olympischer Sportbund (DOSB): \href{https://www.dosb.de/aktuelles/news/detail/auch-der-sport-uebernimmt-verantwortung}{Auch der Sport übernimmt Verantwortung}, 2018.
		\bibitem{sport-prevention-risks} Liechtenstein Olympic Committee: \href{https://www.olympic.li/download_file/view/297/526}{Suchtprävention im Verein: Was können wir tun?} (Präsentationsfolien).
		\bibitem{sport-instead-of-addiction} Bundesministerium für Inneres der Republik Österreich: \href{https://www.bmi.gv.at}{Sport statt Sucht} (Kampagneninformation).
		\bibitem{sport-therapy} AOK Gesundheitsmagazin: \href{https://www.aok.de/pk/magazin/koerper-psyche/sucht/wie-sport-in-der-suchttherapie-helfen-kann/}{Wie Sport in der Suchttherapie helfen kann}, 2025.
		\bibitem{dhs-reitox-2025} Deutsche Hauptstelle für Suchtfragen (DHS): \href{https://www.dhs.de/fileadmin/user_upload/pdf/Unsere_Arbeit/REITOX_BERICHT_2025_DE_WB_04_Praevention.pdf}{REITOX-BERICHT 2025 – Prävention}, S. 2.
		\bibitem{health-literacy-schools} Deutsches Kinderhilfswerk: \href{https://www.dkhw.de/filestorage/1_Informieren/1.1_Unsere_Themen/Kinderrechte/Kinderrechte-Index/2025/Indikatoren_Datengrundlage/Gesundheit/Vermittlung_von_Gesundheitskompetenzen_in_der_Schule_Auswertung.pdf}{Vermittlung von Gesundheitskompetenzen in der Schule}, 2025, S. 1-2.
		% Neue Quellen für Stuttgart und Baden-Württemberg (bereits im vorherigen Code enthalten, ergänzt um Esslingen)
		\bibitem{stuttgart-drug-delikte} Schwäbische Zeitung: \href{https://www.schwaebische.de/regional/baden-wuerttemberg/der-mann-mit-dem-koks-ist-da-so-schnell-liefert-nicht-mal-der-pizzadienst-3125567}{Der Mann mit dem Koks ist da: So schnell liefert nicht mal der Pizzadienst}, 2024.
		\bibitem{sicherheitsbericht-2025} Polizei Baden-Württemberg: \href{https://www.polizei-bw.de/sicherheitsbericht-2025/}{Sicherheitsbericht 2025}.
		\bibitem{drogentote-2025} Innenministerium Baden-Württemberg: \href{https://im.baden-wuerttemberg.de/de/service/presse-und-oeffentlichkeitsarbeit/pressemitteilung/pid/zahl-der-drogentoten-im-jahr-2025-leicht-zurueckgegangen}{Zahl der Drogentoten im Jahr 2025 leicht zurückgegangen}, 2026.
		\bibitem{drogentote-2023} BKZ: \href{https://www.bkz.de/nachrichten/weniger-drogentote-im-suedwesten-erstmals-lachgas-todesfall-234659.html}{Weniger Drogentote im Südwesten - erstmals Lachgas-Todesfall}, 2024.
		\bibitem{drogenaffinitaetsstudie} BZgA/Präventionstag: \href{https://www.praeventionstag.de}{Drogenaffinitätsstudie 2025}, 2025.
		\bibitem{drogenszene-stuttgart} Wikipedia: \href{http://web.archive.org/web/20060913000000/http://de.wikipedia.org:80/wiki/Drogenszene}{Offene Drogenszene in Stuttgart}.
		\bibitem{drogenszene-verlagerung} Deutsche Digitale Bibliothek: \href{https://www.deutsche-digitale-bibliothek.de/item/UC5HIVCN7D2OED4JWS3CDTR4AS7UKAFW}{Offene Drogenszene in Stuttgart}, 1995.
		\bibitem{razzia-stuttgart-sep24} Zeit Online: \href{https://www.zeit.de/news/2024-09/16/durchsuchungen-schlag-gegen-drogenkriminalitaet-im-grossraum-stuttgart}{Durchsuchungen: Schlag gegen Drogenkriminalität im Großraum Stuttgart}, 2024.
		\bibitem{drogennetzwerk-stuttgart} Stimme: \href{https://www.stimme.de/deutschland-welt/polizei-sprengt-drogennetzwerk-im-raum-stuttgart-zehn-personen-festgenommen-51882359.html}{Polizei sprengt Drogennetzwerk im Raum Stuttgart: Zehn Personen festgenommen}, 2024.
		\bibitem{landtag-antrag} Landtag Baden-Württemberg, Drucksache 17/7228: \href{https://www.landtag-bw.de/resource/blob/267420/8434cfbd4572d13f974635de860859b8/17_7228_D.pdf}{Kokstaxis und der Drogenhandel in Baden-Württemberg}, 2024.
		\bibitem{substitution-stuttgart} YouTube: Heroinsubstitution in Stuttgart: Eine Betroffene erzählt.
		\bibitem{suchthilfe-stuttgart} Landeshauptstadt Stuttgart: \href{https://www.stuttgart.de}{Suchthilfe in Stuttgart}.
		\bibitem{suchtberatung-stuttgart} Klinikum Stuttgart: \href{https://www.klinikum-stuttgart.de/sucht}{Suchtberatung Türlenstraße 22}.
		\bibitem{kokainfund-lka} Stuttgarter Zeitung: \href{https://www.stuttgarter-zeitung.de/inhalt.riesiger-kokainfund-tonnenweise-drogen-im-stuttgarter-lka.48cd2e57-7f17-490e-89a8-b0fc6697426c.html}{Größter Kokain-Fund Deutschlands: Tonnenweise Drogen im Stuttgarter LKA gelagert}, 2024.
		% Quellen für Landkreis Esslingen
		\bibitem{esslingen-drogentote} Eßlinger Zeitung: \href{https://www.esslinger-zeitung.de}{Sieben Drogentote im Kreis Esslingen 2023}, 2024.
		\bibitem{esslingen-bande} Polizeipräsidium Reutlingen: Pressemitteilung vom 15. Januar 2025 – Festnahme von sechs mutmaßlichen Bandenmitgliedern.
		\bibitem{esslingen-kokain} Südwest Presse: \href{https://www.swp.de}{Großrazzia gegen Kokainhandel in Kirchheim}, 2024.
		\bibitem{esslingen-plantage} Polizeipräsidium Esslingen: Pressemitteilung vom 10. Oktober 2024 – Cannabisplantage im Stadtgebiet Esslingen ausgehoben.
		\bibitem{esslingen-unfaelle} Landratsamt Esslingen: Verkehrssicherheitsbericht 2025, S. 18-21.
		\bibitem{esslingen-praevention} Landratsamt Esslingen, Suchtprophylaxe: Jahresbericht 2024.
		\bibitem{esslingen-kontakt} Landratsamt Esslingen: \href{https://www.landkreis-esslingen.de/site/LRA-ES-Internet-2023/node/19421567}{Suchtprävention im Landkreis Esslingen}, 2025.
	\end{thebibliography}
\end{document}